\section{Síntesis y ampliación}

\subsection{Panorama general del proceso de Post-Training}

Todo el Post-Training es un ciclo iterativo, formado por tres etapas de igual peso:

\begin{enumerate}
    \item \textbf{Construcción de datos} (aproximadamente 1/3 del tiempo): diseñar la proporción de mezcla de datos, garantizar exactitud, diversidad y complejidad, y producir en lotes datos de entrenamiento de alta calidad mediante un pipeline de seed data + generación con LLM + filtrado
    \item \textbf{Entrenamiento del modelo} (aproximadamente 1/3 del tiempo): SFT (se recomienda LoRA) $\rightarrow$ Preference Alignment (se recomienda DPO) $\rightarrow$ Model Merging opcional
    \item \textbf{Iteración de evaluación} (aproximadamente 1/3 del tiempo): benchmarks automatizados, LLM-as-Judge y evaluación humana en conjunto, con los resultados de evaluación retroalimentando la siguiente ronda de datos y de mejoras en el entrenamiento
\end{enumerate}

\subsection{Repaso de los puntos centrales}

\begin{importantbox}{Takeaways centrales de esta clase}
\begin{itemize}
    \item El Post-Training convierte un base model en un assistant útil; los pasos centrales son SFT + Preference Alignment
    \item La calidad de los datos importa más que la cantidad de datos, en tres dimensiones: Accuracy, Diversity, Complexity
    \item LoRA es la mejor técnica de SFT en la mayoría de los escenarios; DPO es el algoritmo de alineación más práctico
    \item Model Merging es una herramienta poderosa para combinar las ventajas de varios modelos sin costo adicional
    \item La correlación entre la preferencia humana y los benchmarks automatizados es baja, por lo que es necesario combinar varios métodos de evaluación
    \item Test-Time Compute Scaling intercambia tiempo de inferencia por calidad de salida, y es la tendencia más importante en este momento
\end{itemize}
\end{importantbox}

\subsection{Tabla de decisiones de ingeniería para elegir método}

\begin{table}[H]
\centering
\begin{tabular}{p{0.2\textwidth} p{0.24\textwidth} p{0.24\textwidth}}
\toprule
Escenario objetivo & Método prioritario & Criterio de decisión clave \\
\midrule
Validar primero rápidamente la necesidad & In-context learning / RAG & No pasar demasiado pronto al fine-tuning; primero ver si el base model ya es suficientemente bueno \\
Se necesita un formato de salida o un tono estables & SFT (generalmente con LoRA) & Si la calidad de las muestras es suficientemente alta y si el estilo de salida puede demostrarse de manera explícita \\
Se necesita que las respuestas se ajusten más a la preferencia humana & DPO / Preference Alignment & Si se cuenta con datos de preferencia chosen vs rejected, y si la evaluación cubre la preferencia real de los usuarios \\
Se necesita combinar las especialidades de varios modelos & Model Merging & Si las capacidades de los modelos son complementarias, y si existe riesgo de degradación tras la combinación \\
Aún queda presupuesto en la etapa de inferencia & Test-Time Compute Scaling & Si la latencia y el costo permiten intercambiar más muestreo por mayor calidad \\
\bottomrule
\end{tabular}
\caption{Orden de elección de los principales métodos de Post-Training}
\end{table}

\subsection{Orden de implementación para quien lo lleva a la práctica}

Si se convierte esta clase en recomendaciones de ejecución de ingeniería, el orden más seguro suele ser:

\begin{enumerate}
    \item Primero usar prompting / RAG, de costo mínimo, para delimitar el problema
    \item Después usar un conjunto de datos pequeño y de alta calidad para hacer SFT y resolver los problemas de estilo y formato
    \item Si se necesita mejorar aún más la usabilidad, introducir datos de preferencia para hacer DPO
    \item Por último, según el costo, la latencia y el resultado, elegir model merging o test-time scaling como medio de mejora
\end{enumerate}

\subsection{Lecturas adicionales}

\begin{itemize}
    \item \textbf{LLM Engineer's Handbook}: libro de referencia sistemático escrito por el ponente
    \item \textbf{Open Perfect Blend} (Hugging Face): esquema de mezcla de datos de SFT de propósito general
    \item \textbf{mergekit}: herramienta de código abierto para la combinación de modelos
    \item \textbf{TRL} (Hugging Face): la biblioteca de herramientas de Post-Training más completa
    \item Rafailov et al., \textit{Direct Preference Optimization} (NeurIPS 2023): artículo original de DPO
    \item Hu et al., \textit{LoRA: Low-Rank Adaptation of Large Language Models} (ICLR 2022): artículo original de LoRA
    \item Snell et al., \textit{Scaling LLM Test-Time Compute} (2024): fundamento teórico de Test-Time Compute
    \item Chatbot Arena (\url{https://arena.lmsys.org}): clasificación de modelos basada en la preferencia humana
\end{itemize}

%% --- Fin del contenido --- %%

\end{document}
